\documentclass[11pt,a4paper]{article}

\usepackage[margin=2.2cm]{geometry}
\usepackage{kotex}
\usepackage{amsmath,amssymb}
\usepackage{booktabs}
\usepackage{array}
\usepackage{xcolor}
\usepackage{listings}
\usepackage[most]{tcolorbox}
\usepackage{enumitem}
\usepackage{hyperref}
\usepackage{fancyhdr}
\usepackage{titlesec}

\definecolor{codebg}{RGB}{247,247,247}
\definecolor{codeframe}{RGB}{210,210,210}
\definecolor{keywordcolor}{RGB}{0,70,160}
\definecolor{commentcolor}{RGB}{0,120,80}
\definecolor{stringcolor}{RGB}{170,50,50}
\definecolor{titleblue}{RGB}{35,75,130}

\hypersetup{colorlinks=true,linkcolor=titleblue,urlcolor=titleblue}

\pagestyle{fancy}
\fancyhf{}
\lhead{C++ Programming}
\rhead{Day 10}
\cfoot{\thepage}

\titleformat{\section}{\Large\bfseries\color{titleblue}}{\thesection.}{0.6em}{}
\titleformat{\subsection}{\large\bfseries}{\thesubsection.}{0.5em}{}

\lstdefinestyle{cppstyle}{
    language=C++,
    backgroundcolor=\color{codebg},
    frame=single,
    rulecolor=\color{codeframe},
    basicstyle=\ttfamily\small,
    keywordstyle=\color{keywordcolor}\bfseries,
    commentstyle=\color{commentcolor},
    stringstyle=\color{stringcolor},
    numbers=left,
    numberstyle=\tiny\color{gray},
    numbersep=8pt,
    tabsize=4,
    showstringspaces=false,
    breaklines=true,
    columns=fullflexible
}
\lstset{style=cppstyle}

\newtcolorbox{learningbox}{colback=blue!4,colframe=titleblue,title=\textbf{학습 목표},fonttitle=\bfseries}
\newtcolorbox{importantbox}{colback=yellow!8,colframe=orange!70!black,title=\textbf{핵심 포인트},fonttitle=\bfseries}
\newtcolorbox{practicebox}{colback=red!3,colframe=red!55!black,title=\textbf{실습},fonttitle=\bfseries}

\begin{document}

\begin{titlepage}
    \centering
    \vspace*{3cm}
    {\Huge\bfseries C++ Programming\par}
    \vspace{0.8cm}
    {\LARGE Day 10: STL 기초\par}
    \vspace{2cm}
    {\large vector, iterator, pair, algorithm, lambda\par}
    \vfill
\end{titlepage}

\tableofcontents
\newpage

\section{학습 목표}
\begin{learningbox}
이 장을 마치면 다음 내용을 설명하고 직접 코드로 작성할 수 있다.
\begin{itemize}[leftmargin=1.5em]
    \item STL의 의미와 컨테이너, 반복자, 알고리즘의 관계 설명하기
    \item \texttt{vector}를 선언하고 원소를 추가, 접근, 삭제하기
    \item \texttt{size()}, \texttt{empty()}, \texttt{front()}, \texttt{back()} 사용하기
    \item 인덱스 반복문과 범위 기반 반복문으로 \texttt{vector} 순회하기
    \item 반복자의 역할과 \texttt{begin()}, \texttt{end()}의 의미 설명하기
    \item 반복자를 역참조하여 원소를 읽고 수정하기
    \item \texttt{iterator}와 \texttt{const\_iterator}의 차이 구분하기
    \item \texttt{pair}로 서로 관련된 두 값을 하나의 객체에 저장하기
    \item \texttt{first}, \texttt{second}, \texttt{make\_pair()} 사용하기
    \item \texttt{sort()}, \texttt{find()}, \texttt{count()}, \texttt{reverse()} 사용하기
    \item 알고리즘 함수에 반복자 범위를 전달하는 방식 이해하기
    \item 람다 표현식의 기본 문법과 매개변수, 반환값 이해하기
    \item 람다를 \texttt{sort()}와 함께 사용하여 정렬 기준 정의하기
    \item 캡처 목록을 사용하여 외부 변수를 읽거나 수정하기
\end{itemize}
\end{learningbox}

\section{STL의 개요}
STL(Standard Template Library)은 자주 사용하는 자료구조와 알고리즘을 템플릿으로 제공하는 C++ 표준 라이브러리의 핵심 부분이다. Day 9에서 학습한 클래스 템플릿과 함수 템플릿이 실제 라이브러리에서 어떻게 활용되는지 보여 주는 대표적인 예가 STL이다.

STL은 크게 컨테이너(container), 반복자(iterator), 알고리즘(algorithm)으로 이해할 수 있다.

\begin{center}
\begin{tabular}{lll}
\toprule
구성 요소 & 역할 & 예시 \\
\midrule
컨테이너 & 여러 데이터를 저장하는 자료구조 & \texttt{vector}, \texttt{list}, \texttt{map} \\
반복자 & 컨테이너의 원소 위치를 가리키는 객체 & \texttt{begin()}, \texttt{end()} \\
알고리즘 & 반복자 범위의 데이터를 처리하는 함수 & \texttt{sort()}, \texttt{find()} \\
\bottomrule
\end{tabular}
\end{center}

예를 들어 \texttt{vector<int>}는 정수를 저장하는 컨테이너이고, \texttt{numbers.begin()}과 \texttt{numbers.end()}는 원소 범위를 나타내는 반복자를 제공한다. \texttt{sort()}는 이 반복자 범위를 받아 실제 정렬을 수행한다.

\begin{lstlisting}
vector<int> numbers = {30, 10, 20};
sort(numbers.begin(), numbers.end());
\end{lstlisting}

위 문장은 다음 구조로 해석할 수 있다.

\begin{verbatim}
Container: vector<int> numbers
Range:     [numbers.begin(), numbers.end())
Algorithm: sort
\end{verbatim}

범위의 마지막에 있는 \texttt{end()}는 마지막 원소 그 자체가 아니라 마지막 원소 다음 위치를 나타낸다. 따라서 STL의 기본 범위는 시작 위치를 포함하고 끝 위치는 포함하지 않는 반개구간 형태이다.

\begin{verbatim}
[begin, end)
 begin is included
 end is excluded
\end{verbatim}

\begin{importantbox}
STL에서는 컨테이너가 데이터를 저장하고, 반복자가 데이터의 범위를 표현하며, 알고리즘이 그 범위를 처리한다. 이 세 요소가 서로 분리되어 있기 때문에 같은 알고리즘을 여러 컨테이너에 재사용할 수 있다.
\end{importantbox}

\section{vector}
\texttt{vector}는 같은 자료형의 여러 값을 연속적으로 저장하는 동적 배열 컨테이너이다. 일반 배열과 달리 실행 중에 원소를 추가하거나 삭제하면서 크기를 변경할 수 있다.

\subsection{헤더와 기본 선언}
\texttt{vector}를 사용하려면 \texttt{<vector>} 헤더를 포함한다.

\begin{lstlisting}
#include <vector>
\end{lstlisting}

\texttt{vector}는 클래스 템플릿이므로 각괄호 안에 저장할 원소의 자료형을 지정한다.

\begin{lstlisting}
vector<int> numbers;
vector<double> values;
vector<string> names;
\end{lstlisting}

각각 정수, 실수, 문자열을 저장하는 서로 다른 \texttt{vector} 타입이다.

\subsection{초기화}
빈 \texttt{vector}를 만든 뒤 원소를 추가할 수도 있고, 선언과 동시에 초기값을 지정할 수도 있다.

\begin{lstlisting}
vector<int> numbers;
vector<int> scores = {90, 85, 100};
\end{lstlisting}

같은 값으로 여러 원소를 초기화하는 문법도 있다.

\begin{lstlisting}
vector<int> values(5, 10);
\end{lstlisting}

위 코드는 값이 모두 \texttt{10}인 정수 원소 다섯 개를 만든다.

\begin{verbatim}
10 10 10 10 10
\end{verbatim}

다음 코드는 원소 다섯 개를 기본값으로 초기화한다.

\begin{lstlisting}
vector<int> values(5);
\end{lstlisting}

정수의 기본값은 \texttt{0}이므로 다섯 원소가 모두 \texttt{0}이 된다.

\subsection{원소 추가와 삭제}
\texttt{push\_back()}은 \texttt{vector}의 끝에 새 원소를 추가한다.

\begin{lstlisting}
vector<string> fruits;

fruits.push_back("Apple");
fruits.push_back("Banana");
fruits.push_back("Orange");
\end{lstlisting}

\texttt{pop\_back()}은 마지막 원소를 제거한다.

\begin{lstlisting}
fruits.pop_back();
\end{lstlisting}

\texttt{pop\_back()}은 제거한 값을 반환하지 않는다. 또한 빈 \texttt{vector}에서 호출하면 올바른 동작이 보장되지 않으므로 필요한 경우 \texttt{empty()}로 먼저 확인한다.

\begin{lstlisting}
if (!fruits.empty())
{
    fruits.pop_back();
}
\end{lstlisting}

\subsection{크기와 비어 있는지 확인}
\texttt{size()}는 현재 저장된 원소의 개수를 반환한다.

\begin{lstlisting}
cout << fruits.size() << endl;
\end{lstlisting}

\texttt{empty()}는 원소가 하나도 없으면 \texttt{true}, 하나 이상 있으면 \texttt{false}를 반환한다.

\begin{lstlisting}
if (fruits.empty())
{
    cout << "Empty" << endl;
}
\end{lstlisting}

\subsection{인덱스로 원소 접근}
배열처럼 대괄호를 사용하여 원소에 접근할 수 있다.

\begin{lstlisting}
cout << fruits[0] << endl;
fruits[1] = "Mango";
\end{lstlisting}

\texttt{operator[]}는 인덱스의 유효 범위를 검사하지 않는다. 잘못된 인덱스를 사용하면 정의되지 않은 동작이 발생할 수 있다.

\texttt{at()}은 범위를 검사하는 접근 함수이다.

\begin{lstlisting}
cout << fruits.at(0) << endl;
\end{lstlisting}

잘못된 인덱스를 전달하면 예외가 발생하므로 오류를 더 명확하게 발견할 수 있다.

\subsection{첫 번째와 마지막 원소}
\texttt{front()}는 첫 번째 원소, \texttt{back()}은 마지막 원소에 접근한다.

\begin{lstlisting}
cout << fruits.front() << endl;
cout << fruits.back() << endl;
\end{lstlisting}

두 함수 모두 빈 \texttt{vector}에 사용하면 안 된다.

\subsection{인덱스 반복문}
\texttt{size()}를 사용하여 모든 원소를 순회할 수 있다.

\begin{lstlisting}
for (size_t i = 0; i < fruits.size(); i++)
{
    cout << fruits[i] << endl;
}
\end{lstlisting}

\texttt{size()}의 반환형은 부호 없는 정수 계열이므로 인덱스 변수에 \texttt{size\_t}를 사용하면 자료형이 자연스럽게 맞는다.

\subsection{범위 기반 반복문}
원소의 인덱스가 필요하지 않다면 범위 기반 반복문을 사용할 수 있다.

\begin{lstlisting}
for (string fruit : fruits)
{
    cout << fruit << endl;
}
\end{lstlisting}

위 코드는 각 원소를 \texttt{fruit}에 복사한다. 원소를 직접 수정하려면 참조를 사용한다.

\begin{lstlisting}
for (string& fruit : fruits)
{
    fruit += " Juice";
}
\end{lstlisting}

읽기만 하면서 불필요한 복사를 피하려면 상수 참조를 사용할 수 있다.

\begin{lstlisting}
for (const string& fruit : fruits)
{
    cout << fruit << endl;
}
\end{lstlisting}

\subsection{전체 예제}
\begin{lstlisting}
#include <iostream>
#include <vector>
#include <string>

using namespace std;

int main()
{
    vector<string> fruits;

    fruits.push_back("Apple");
    fruits.push_back("Banana");
    fruits.push_back("Orange");
    fruits.push_back("Grape");

    cout << "Size: " << fruits.size() << endl;
    cout << "First: " << fruits.front() << endl;
    cout << "Last: " << fruits.back() << endl;

    for (size_t i = 0; i < fruits.size(); i++)
    {
        cout << fruits[i] << endl;
    }

    fruits.pop_back();

    cout << "After pop_back" << endl;

    for (const string& fruit : fruits)
    {
        cout << fruit << endl;
    }

    return 0;
}
\end{lstlisting}

출력 결과는 다음과 같다.

\begin{verbatim}
Size: 4
First: Apple
Last: Grape
Apple
Banana
Orange
Grape
After pop_back
Apple
Banana
Orange
\end{verbatim}

\begin{importantbox}
\texttt{vector}는 크기가 변할 수 있는 동적 배열이다. 끝에 원소를 추가할 때는 \texttt{push\_back()}, 마지막 원소를 제거할 때는 \texttt{pop\_back()}, 현재 원소 수를 확인할 때는 \texttt{size()}를 사용한다.
\end{importantbox}

\section{iterator}
반복자(iterator)는 컨테이너의 특정 원소 위치를 가리키는 객체이다. 포인터와 비슷하게 이동할 수 있고, 역참조 연산자 \texttt{*}를 사용하여 해당 위치의 원소에 접근할 수 있다.

\subsection{begin과 end}
\texttt{begin()}은 첫 번째 원소의 위치를 나타내는 반복자를 반환한다. \texttt{end()}는 마지막 원소 다음 위치를 나타내는 반복자를 반환한다.

\begin{verbatim}
vector:   [5] [10] [15] [20]
           ^                 ^
         begin              end
\end{verbatim}

\texttt{end()}가 가리키는 위치에는 실제 원소가 없으므로 역참조하면 안 된다.

\subsection{반복자 선언}
\texttt{vector<int>}의 반복자 타입은 다음과 같이 작성한다.

\begin{lstlisting}
vector<int>::iterator it;
\end{lstlisting}

반복자의 전체 타입이 길어질 수 있으므로 \texttt{auto}를 자주 사용한다.

\begin{lstlisting}
auto it = numbers.begin();
\end{lstlisting}

컴파일러가 \texttt{numbers.begin()}의 반환형으로부터 반복자 타입을 추론한다.

\subsection{반복자를 이용한 순회}
첫 번째 원소에서 시작하여 \texttt{end()}에 도달할 때까지 반복자를 증가시킨다.

\begin{lstlisting}
for (auto it = numbers.begin(); it != numbers.end(); it++)
{
    cout << *it << " ";
}
\end{lstlisting}

각 부분의 의미는 다음과 같다.

\begin{center}
\begin{tabular}{ll}
\toprule
표현 & 의미 \\
\midrule
\texttt{numbers.begin()} & 첫 번째 원소 위치 \\
\texttt{numbers.end()} & 마지막 원소 다음 위치 \\
\texttt{it != numbers.end()} & 아직 범위 안에 있는지 확인 \\
\texttt{it++} & 다음 원소 위치로 이동 \\
\texttt{*it} & 현재 위치의 원소에 접근 \\
\bottomrule
\end{tabular}
\end{center}

\subsection{반복자로 원소 수정}
일반 반복자를 역참조한 결과는 실제 원소이므로 값을 수정할 수 있다.

\begin{lstlisting}
for (auto it = numbers.begin(); it != numbers.end(); it++)
{
    *it += 3;
}
\end{lstlisting}

각 원소에 \texttt{3}이 더해진다.

\subsection{const\_iterator}
원소를 읽기만 하고 수정하지 못하게 하려면 상수 반복자를 사용한다.

\begin{lstlisting}
vector<int>::const_iterator it;
\end{lstlisting}

\texttt{cbegin()}과 \texttt{cend()}는 상수 반복자를 반환한다.

\begin{lstlisting}
for (auto it = numbers.cbegin(); it != numbers.cend(); it++)
{
    cout << *it << " ";
}
\end{lstlisting}

다음과 같이 값을 수정하려고 하면 컴파일 오류가 발생한다.

\begin{lstlisting}
for (auto it = numbers.cbegin(); it != numbers.cend(); it++)
{
    *it = 100;
}
\end{lstlisting}

\texttt{*it}가 읽기 전용으로 취급되기 때문이다.

\subsection{반복자와 포인터의 유사성}
반복자는 포인터와 완전히 같은 타입은 아니지만 사용 방식이 비슷하다.

\begin{center}
\begin{tabular}{lll}
\toprule
동작 & 포인터 & 반복자 \\
\midrule
현재 값 접근 & \texttt{*ptr} & \texttt{*it} \\
다음 위치 이동 & \texttt{ptr++} & \texttt{it++} \\
위치 비교 & \texttt{ptr != endPtr} & \texttt{it != container.end()} \\
\bottomrule
\end{tabular}
\end{center}

반복자는 컨테이너의 구현 방식에 맞는 공통 접근 방법을 제공한다. 알고리즘 함수가 컨테이너 자체가 아니라 반복자 범위를 받는 이유도 여기에 있다.

\subsection{반복자 무효화}
\texttt{vector}에 원소를 추가하거나 삭제하면 내부 저장 공간이 재배치될 수 있다. 이 경우 기존 반복자가 더 이상 유효하지 않을 수 있다.

\begin{lstlisting}
auto it = numbers.begin();
numbers.push_back(100);
\end{lstlisting}

\texttt{push\_back()} 과정에서 재할당이 발생했다면 \texttt{it}를 계속 사용하면 안 된다. 컨테이너를 변경한 뒤에는 필요한 반복자를 다시 얻는 습관이 안전하다.

\subsection{전체 예제}
\begin{lstlisting}
#include <iostream>
#include <vector>

using namespace std;

int main()
{
    vector<int> numbers = {5, 10, 15, 20, 25};

    for (auto it = numbers.begin(); it != numbers.end(); it++)
    {
        cout << *it << " ";
    }

    cout << endl;

    for (auto it = numbers.begin(); it != numbers.end(); it++)
    {
        *it += 3;
    }

    for (auto it = numbers.cbegin(); it != numbers.cend(); it++)
    {
        cout << *it << " ";
    }

    cout << endl;

    return 0;
}
\end{lstlisting}

출력 결과는 다음과 같다.

\begin{verbatim}
5 10 15 20 25
8 13 18 23 28
\end{verbatim}

\begin{importantbox}
반복자는 컨테이너의 원소 위치를 나타낸다. \texttt{begin()}부터 시작하여 \texttt{end()}에 도달하기 전까지 이동하고, \texttt{*it}로 현재 원소에 접근한다. 읽기 전용 순회에는 \texttt{cbegin()}과 \texttt{cend()}를 사용할 수 있다.
\end{importantbox}

\section{pair}
\texttt{pair}는 서로 관련된 두 값을 하나의 객체에 함께 저장하는 클래스 템플릿이다. 두 값은 서로 같은 자료형일 수도 있고 서로 다른 자료형일 수도 있다.

\subsection{헤더와 기본 선언}
\texttt{pair}는 \texttt{<utility>} 헤더에 정의되어 있다.

\begin{lstlisting}
#include <utility>
\end{lstlisting}

두 타입을 각괄호 안에 순서대로 지정한다.

\begin{lstlisting}
pair<string, int> student;
\end{lstlisting}

위 객체는 첫 번째 값으로 \texttt{string}, 두 번째 값으로 \texttt{int}를 저장한다.

\subsection{초기화}
중괄호를 사용하여 선언과 동시에 두 값을 초기화할 수 있다.

\begin{lstlisting}
pair<string, int> student = {"Alice", 95};
\end{lstlisting}

생성자 형태로도 초기화할 수 있다.

\begin{lstlisting}
pair<string, int> student("Alice", 95);
\end{lstlisting}

\subsection{first와 second}
첫 번째 값은 \texttt{first}, 두 번째 값은 \texttt{second} 멤버로 접근한다.

\begin{lstlisting}
cout << student.first << endl;
cout << student.second << endl;
\end{lstlisting}

두 멤버는 읽는 것뿐만 아니라 수정할 수도 있다.

\begin{lstlisting}
student.first = "Bob";
student.second = 88;
\end{lstlisting}

\subsection{make\_pair}
\texttt{make\_pair()}를 사용하면 전달된 값으로부터 두 자료형을 추론하여 \texttt{pair} 객체를 생성할 수 있다.

\begin{lstlisting}
auto student = make_pair(string("Alice"), 95);
\end{lstlisting}

문자열 리터럴을 그대로 사용하면 첫 번째 타입이 문자열 포인터 계열로 추론될 수 있다. 명확한 \texttt{pair<string, int>}를 원한다면 \texttt{string("Alice")}처럼 \texttt{string} 객체를 전달하거나 타입을 직접 선언한다.

\begin{lstlisting}
pair<string, int> student = make_pair("Alice", 95);
\end{lstlisting}

대입 과정에서 첫 번째 값이 \texttt{string}으로 변환된다.

\subsection{pair 비교}
같은 타입의 \texttt{pair}끼리는 비교할 수 있다. 먼저 \texttt{first}를 비교하고, \texttt{first}가 같으면 \texttt{second}를 비교한다.

\begin{lstlisting}
pair<int, int> a = {1, 20};
pair<int, int> b = {2, 10};

cout << (a < b) << endl;
\end{lstlisting}

첫 번째 값에서 \texttt{1 < 2}이므로 결과는 \texttt{true}이다. 이 비교 규칙은 \texttt{vector<pair<...>>}를 기본 정렬할 때도 적용된다.

\subsection{vector와 pair 함께 사용하기}
여러 학생의 이름과 점수를 저장하려면 \texttt{pair<string, int>}를 원소 타입으로 하는 \texttt{vector}를 만들 수 있다.

\begin{lstlisting}
vector<pair<string, int>> students;

students.push_back({"Alice", 95});
students.push_back({"Bob", 88});
students.push_back({"Charlie", 92});
\end{lstlisting}

각 원소는 하나의 \texttt{pair} 객체이다.

\begin{lstlisting}
for (const auto& student : students)
{
    cout << student.first << " " << student.second << endl;
}
\end{lstlisting}

\texttt{auto}를 사용하면 긴 타입 이름을 반복해서 작성하지 않아도 된다. \texttt{const auto\&}는 복사 없이 읽기 전용으로 각 원소에 접근한다.

\subsection{전체 예제}
\begin{lstlisting}
#include <iostream>
#include <vector>
#include <string>
#include <utility>

using namespace std;

int main()
{
    pair<string, int> student = {"Alice", 95};

    cout << student.first << endl;
    cout << student.second << endl;

    student.second = 98;

    vector<pair<string, int>> students;

    students.push_back(student);
    students.push_back(make_pair(string("Bob"), 88));
    students.push_back({"Charlie", 92});

    for (const auto& item : students)
    {
        cout << item.first << ": " << item.second << endl;
    }

    return 0;
}
\end{lstlisting}

출력 결과는 다음과 같다.

\begin{verbatim}
Alice
95
Alice: 98
Bob: 88
Charlie: 92
\end{verbatim}

\begin{importantbox}
\texttt{pair<T, U>}는 두 값을 하나로 묶는다. 첫 번째 값은 \texttt{first}, 두 번째 값은 \texttt{second}로 접근한다. \texttt{vector<pair<T, U>>} 형태를 사용하면 서로 관련된 두 값의 목록을 쉽게 관리할 수 있다.
\end{importantbox}

\section{algorithm}
\texttt{<algorithm>} 헤더는 정렬, 탐색, 개수 계산, 순서 뒤집기 등 자주 사용하는 알고리즘 함수를 제공한다. 대부분의 알고리즘은 처리할 컨테이너 자체가 아니라 반복자로 표현한 범위를 전달받는다.

\subsection{헤더 포함}
\begin{lstlisting}
#include <algorithm>
\end{lstlisting}

\subsection{sort}
\texttt{sort()}는 지정한 반복자 범위의 원소를 기본적으로 오름차순 정렬한다.

\begin{lstlisting}
vector<int> numbers = {40, 10, 30, 20};

sort(numbers.begin(), numbers.end());
\end{lstlisting}

정렬 결과는 다음과 같다.

\begin{verbatim}
10 20 30 40
\end{verbatim}

내림차순 정렬에는 \texttt{greater<int>()}를 사용할 수 있다. 이 기능을 사용할 때는 일반적으로 \texttt{<functional>} 헤더를 포함한다.

\begin{lstlisting}
sort(numbers.begin(), numbers.end(), greater<int>());
\end{lstlisting}

정렬 결과는 다음과 같다.

\begin{verbatim}
40 30 20 10
\end{verbatim}

\subsection{find}
\texttt{find()}는 범위에서 지정한 값과 같은 첫 번째 원소를 찾는다.

\begin{lstlisting}
auto it = find(numbers.begin(), numbers.end(), 30);
\end{lstlisting}

값을 찾으면 해당 원소를 가리키는 반복자를 반환하고, 찾지 못하면 두 번째 인수로 전달한 끝 반복자를 반환한다.

\begin{lstlisting}
if (it != numbers.end())
{
    cout << "Found: " << *it << endl;
}
else
{
    cout << "Not found" << endl;
}
\end{lstlisting}

반환값을 바로 역참조하기 전에 반드시 \texttt{end()}와 비교해야 한다. 찾지 못한 경우 반환된 \texttt{end()}를 역참조하면 안 된다.

\subsection{count}
\texttt{count()}는 범위에서 특정 값이 등장하는 횟수를 반환한다.

\begin{lstlisting}
vector<int> values = {1, 2, 2, 3, 2, 4};

int result = count(values.begin(), values.end(), 2);
\end{lstlisting}

\texttt{result}에는 \texttt{3}이 저장된다.

\subsection{reverse}
\texttt{reverse()}는 지정한 범위의 원소 순서를 반대로 뒤집는다.

\begin{lstlisting}
reverse(values.begin(), values.end());
\end{lstlisting}

\texttt{sort()}와 달리 크기 순서로 정렬하는 것이 아니라 현재 순서를 그대로 반전한다.

\subsection{일부 범위만 처리하기}
STL 알고리즘은 \texttt{begin()}과 \texttt{end()} 전체 범위뿐 아니라 일부 범위에도 적용할 수 있다.

\begin{lstlisting}
vector<int> numbers = {50, 40, 30, 20, 10};

sort(numbers.begin(), numbers.begin() + 3);
\end{lstlisting}

처리 범위는 \texttt{[numbers.begin(), numbers.begin() + 3)}이므로 앞의 세 원소만 정렬된다.

\begin{verbatim}
30 40 50 20 10
\end{verbatim}

\texttt{vector}의 반복자는 임의 접근 반복자이므로 \texttt{begin() + n} 형태의 이동이 가능하다.

\subsection{pair의 기본 정렬}
\texttt{vector<pair<int, string>>}를 기본 \texttt{sort()}로 정렬하면 \texttt{pair}의 비교 규칙이 사용된다. 먼저 \texttt{first}를 기준으로 정렬하고, \texttt{first}가 같으면 \texttt{second}를 기준으로 정렬한다.

\begin{lstlisting}
vector<pair<int, string>> items = {
    {2, "Banana"},
    {1, "Orange"},
    {1, "Apple"}
};

sort(items.begin(), items.end());
\end{lstlisting}

정렬 결과는 다음 순서가 된다.

\begin{verbatim}
1 Apple
1 Orange
2 Banana
\end{verbatim}

\subsection{전체 예제}
\begin{lstlisting}
#include <iostream>
#include <vector>
#include <algorithm>

using namespace std;

int main()
{
    vector<int> numbers = {40, 10, 30, 20, 30};

    sort(numbers.begin(), numbers.end());

    for (int number : numbers)
    {
        cout << number << " ";
    }

    cout << endl;

    auto it = find(numbers.begin(), numbers.end(), 30);

    if (it != numbers.end())
    {
        cout << "Found: " << *it << endl;
    }
    else
    {
        cout << "Not found" << endl;
    }

    cout << "Count: "
         << count(numbers.begin(), numbers.end(), 30)
         << endl;

    reverse(numbers.begin(), numbers.end());

    for (int number : numbers)
    {
        cout << number << " ";
    }

    cout << endl;

    return 0;
}
\end{lstlisting}

출력 결과는 다음과 같다.

\begin{verbatim}
10 20 30 30 40
Found: 30
Count: 2
40 30 30 20 10
\end{verbatim}

\begin{importantbox}
STL 알고리즘은 보통 \texttt{[first, last)} 형태의 반복자 범위를 받는다. \texttt{sort()}는 정렬하고, \texttt{find()}는 첫 번째 일치 원소의 반복자를 반환하며, \texttt{count()}는 일치 횟수를 반환하고, \texttt{reverse()}는 현재 순서를 뒤집는다.
\end{importantbox}

\section{lambda}
람다 표현식(lambda expression)은 이름 없는 함수를 코드 안에서 직접 만드는 문법이다. 한 번만 사용할 짧은 동작이나 STL 알고리즘에 전달할 정렬 기준을 작성할 때 특히 유용하다.

\subsection{기본 문법}
람다 표현식의 기본 구조는 다음과 같다.

\begin{lstlisting}
[capture](parameters)
{
    statements
}
\end{lstlisting}

각 부분의 역할은 다음과 같다.

\begin{center}
\begin{tabular}{ll}
\toprule
부분 & 의미 \\
\midrule
\texttt{[capture]} & 외부 변수를 람다 내부에서 사용하는 방법 \\
\texttt{(parameters)} & 람다가 전달받는 매개변수 \\
\texttt{\{ statements \}} & 실행할 코드 \\
\bottomrule
\end{tabular}
\end{center}

가장 간단한 람다는 다음과 같다.

\begin{lstlisting}
[]()
{
    cout << "Hello" << endl;
};
\end{lstlisting}

그러나 위 람다는 정의만 하고 호출하지 않았다. 즉시 호출하려면 뒤에 괄호를 붙인다.

\begin{lstlisting}
[]()
{
    cout << "Hello" << endl;
}();
\end{lstlisting}

\subsection{변수에 저장하기}
람다 객체를 \texttt{auto} 변수에 저장한 뒤 일반 함수처럼 호출할 수 있다.

\begin{lstlisting}
auto greet = []()
{
    cout << "Hello" << endl;
};

greet();
\end{lstlisting}

람다의 정확한 타입은 컴파일러가 생성하므로 보통 \texttt{auto}를 사용한다.

\subsection{매개변수와 반환값}
람다는 일반 함수처럼 매개변수를 받을 수 있다.

\begin{lstlisting}
auto add = [](int a, int b)
{
    return a + b;
};

cout << add(10, 20) << endl;
\end{lstlisting}

반환형은 \texttt{return} 문으로부터 자동 추론된다. 필요하다면 후행 반환형 문법으로 직접 지정할 수 있다.

\begin{lstlisting}
auto divide = [](double a, double b) -> double
{
    return a / b;
};
\end{lstlisting}

\subsection{값 캡처}
람다 바깥의 지역 변수를 사용하려면 캡처 목록에 지정해야 한다.

\begin{lstlisting}
int bonus = 5;

auto addBonus = [bonus](int score)
{
    return score + bonus;
};
\end{lstlisting}

\texttt{[bonus]}는 \texttt{bonus}를 값으로 복사하여 람다 내부에 저장한다. 이후 외부의 \texttt{bonus}가 바뀌어도 이미 캡처된 값에는 영향을 주지 않는다.

\begin{lstlisting}
int bonus = 5;

auto addBonus = [bonus](int score)
{
    return score + bonus;
};

bonus = 20;
cout << addBonus(90) << endl;
\end{lstlisting}

출력은 \texttt{95}이다. 람다를 만들 때의 값 \texttt{5}가 복사되었기 때문이다.

\subsection{참조 캡처}
외부 변수를 직접 읽고 수정하려면 참조로 캡처한다.

\begin{lstlisting}
int total = 0;

auto addToTotal = [&total](int value)
{
    total += value;
};

addToTotal(10);
addToTotal(20);
\end{lstlisting}

람다 실행 후 외부 변수 \texttt{total}은 \texttt{30}이 된다.

\subsection{기본 캡처 방식}
\texttt{[=]}는 사용한 외부 변수를 기본적으로 값으로 캡처한다.

\begin{lstlisting}
auto function = [=]()
{
    cout << a << " " << b << endl;
};
\end{lstlisting}

\texttt{[\&]}는 사용한 외부 변수를 기본적으로 참조로 캡처한다.

\begin{lstlisting}
auto function = [&]()
{
    a++;
    b++;
};
\end{lstlisting}

짧은 예제에서는 편리하지만, 어떤 변수가 어떤 방식으로 캡처되는지 불분명해질 수 있으므로 필요한 변수만 명시적으로 적는 방식이 더 읽기 쉬운 경우가 많다.

\subsection{sort와 람다}
\texttt{sort()}의 세 번째 인수에는 두 원소의 순서를 결정하는 비교 함수를 전달할 수 있다. 람다는 이 비교 기준을 코드 가까이에 작성하게 한다.

\begin{lstlisting}
vector<int> numbers = {40, 10, 30, 20};

sort(numbers.begin(), numbers.end(), [](int a, int b)
{
    return a > b;
});
\end{lstlisting}

비교 람다가 \texttt{a > b}를 반환하므로 큰 값이 앞에 오는 내림차순으로 정렬된다.

비교 함수는 첫 번째 원소가 두 번째 원소보다 앞에 와야 할 때 \texttt{true}를 반환하도록 작성한다.

\begin{verbatim}
return a < b;   ascending order
return a > b;   descending order
\end{verbatim}

\subsection{pair를 원하는 기준으로 정렬하기}
학생 이름과 점수를 저장한 \texttt{pair<string, int>} 목록을 점수 기준 내림차순으로 정렬하려면 \texttt{second}를 비교한다.

\begin{lstlisting}
sort(students.begin(), students.end(),
    [](const pair<string, int>& a,
       const pair<string, int>& b)
    {
        return a.second > b.second;
    });
\end{lstlisting}

두 점수가 같을 때 이름 오름차순까지 적용하려면 조건을 추가한다.

\begin{lstlisting}
sort(students.begin(), students.end(),
    [](const pair<string, int>& a,
       const pair<string, int>& b)
    {
        if (a.second != b.second)
        {
            return a.second > b.second;
        }

        return a.first < b.first;
    });
\end{lstlisting}

\texttt{pair}를 상수 참조로 받으면 비교 과정에서 객체를 복사하지 않고, 람다 내부에서 원소를 수정하지도 못하게 한다.

\subsection{전체 예제}
\begin{lstlisting}
#include <iostream>
#include <vector>
#include <string>
#include <utility>
#include <algorithm>

using namespace std;

int main()
{
    vector<pair<string, int>> students = {
        {"Alice", 92},
        {"Bob", 85},
        {"Charlie", 92},
        {"David", 78}
    };

    sort(students.begin(), students.end(),
        [](const pair<string, int>& a,
           const pair<string, int>& b)
        {
            if (a.second != b.second)
            {
                return a.second > b.second;
            }

            return a.first < b.first;
        });

    for (const auto& student : students)
    {
        cout << student.first
             << ": "
             << student.second
             << endl;
    }

    int bonus = 5;

    auto addBonus = [bonus](int score)
    {
        return score + bonus;
    };

    cout << "Alice with bonus: "
         << addBonus(92)
         << endl;

    return 0;
}
\end{lstlisting}

출력 결과는 다음과 같다.

\begin{verbatim}
Alice: 92
Charlie: 92
Bob: 85
David: 78
Alice with bonus: 97
\end{verbatim}

\begin{importantbox}
람다는 이름 없는 함수 객체이다. \texttt{[]}에는 외부 변수의 캡처 방법을 적고, \texttt{()}에는 매개변수를 적는다. STL 알고리즘의 비교 기준처럼 짧고 특정 위치에서만 필요한 동작을 작성할 때 유용하다.
\end{importantbox}

\section{STL 요소의 연결}
Day 10에서 학습한 다섯 요소는 서로 독립된 문법이 아니라 하나의 처리 흐름으로 연결된다.

\begin{verbatim}
vector<pair<string, int>>
        |
        | stores data
        v
begin() and end()
        |
        | define a range
        v
sort()
        |
        | uses a comparison rule
        v
lambda
\end{verbatim}

다음 코드를 단계별로 해석해 보자.

\begin{lstlisting}
vector<pair<string, int>> students = {
    {"Alice", 92},
    {"Bob", 85},
    {"Charlie", 95}
};

sort(students.begin(), students.end(),
    [](const auto& a, const auto& b)
    {
        return a.second > b.second;
    });
\end{lstlisting}

\begin{enumerate}[leftmargin=2em]
    \item \texttt{vector}가 여러 학생 정보를 저장한다.
    \item 각 학생 정보는 이름과 점수를 묶은 \texttt{pair}이다.
    \item \texttt{begin()}과 \texttt{end()}가 정렬할 전체 범위를 만든다.
    \item \texttt{sort()}가 범위의 원소를 재배치한다.
    \item 람다가 두 \texttt{pair}의 \texttt{second}를 비교하여 점수 내림차순 기준을 제공한다.
\end{enumerate}

이 구조는 이후 더 다양한 STL 컨테이너와 알고리즘을 배울 때도 반복된다.

\begin{importantbox}
STL 코드를 읽을 때는 ``무엇을 저장하는가'', ``어떤 범위를 처리하는가'', ``어떤 알고리즘을 적용하는가'', ``비교 또는 처리 기준은 무엇인가''의 순서로 나누어 해석하면 이해하기 쉽다.
\end{importantbox}

\section{종합 정리}
\begin{importantbox}
\begin{itemize}[leftmargin=1.5em]
    \item STL은 컨테이너, 반복자, 알고리즘을 중심으로 구성된다.
    \item \texttt{vector<T>}는 같은 타입의 여러 원소를 저장하는 동적 배열이다.
    \item \texttt{push\_back()}은 끝에 원소를 추가하고 \texttt{pop\_back()}은 마지막 원소를 제거한다.
    \item \texttt{size()}는 원소 수를 반환하고 \texttt{empty()}는 비어 있는지 확인한다.
    \item \texttt{front()}와 \texttt{back()}은 첫 번째와 마지막 원소에 접근한다.
    \item 범위 기반 반복문에서 참조를 사용하면 실제 원소를 수정할 수 있다.
    \item 반복자는 컨테이너의 원소 위치를 나타내며 포인터와 유사하게 사용한다.
    \item \texttt{begin()}은 첫 번째 원소, \texttt{end()}는 마지막 원소 다음 위치를 나타낸다.
    \item \texttt{*it}는 현재 반복자가 가리키는 원소에 접근한다.
    \item \texttt{cbegin()}과 \texttt{cend()}는 읽기 전용 반복자를 제공한다.
    \item \texttt{pair<T, U>}는 서로 관련된 두 값을 하나로 묶는다.
    \item \texttt{pair}의 값은 \texttt{first}와 \texttt{second}로 접근한다.
    \item \texttt{sort()}, \texttt{find()}, \texttt{count()}, \texttt{reverse()}는 \texttt{<algorithm>}에 정의되어 있다.
    \item STL 알고리즘은 일반적으로 \texttt{[first, last)} 반복자 범위를 처리한다.
    \item \texttt{find()}가 값을 찾지 못하면 끝 반복자를 반환하므로 역참조 전에 비교해야 한다.
    \item 람다는 이름 없는 함수 객체이며 \texttt{[capture](parameters) \{ body \}} 형태로 작성한다.
    \item 값 캡처는 외부 값을 복사하고, 참조 캡처는 외부 변수 자체에 접근한다.
    \item \texttt{sort()}에 람다를 전달하면 사용자 정의 정렬 기준을 작성할 수 있다.
    \item \texttt{vector}, \texttt{iterator}, \texttt{pair}, \texttt{algorithm}, \texttt{lambda}는 실제 STL 코드에서 함께 연결되어 사용된다.
\end{itemize}
\end{importantbox}

\section{개념별 실습 완성 코드}
\subsection{vector}
\begin{practicebox}
문자열을 저장하는 \texttt{vector}를 만들고 과일 네 개를 추가한다. 원소 수와 모든 원소를 출력한 뒤 마지막 원소를 제거하고 다시 출력한다.
\end{practicebox}

\begin{lstlisting}
#include <iostream>
#include <vector>
#include <string>

using namespace std;

int main()
{
    vector<string> fruits;

    fruits.push_back("Apple");
    fruits.push_back("Banana");
    fruits.push_back("Orange");
    fruits.push_back("Grape");

    cout << fruits.size() << endl;

    for (size_t i = 0; i < fruits.size(); i++)
    {
        cout << fruits[i] << endl;
    }

    fruits.pop_back();

    cout << "After pop_back" << endl;

    for (const string& fruit : fruits)
    {
        cout << fruit << endl;
    }

    return 0;
}
\end{lstlisting}

\subsection{iterator}
\begin{practicebox}
정수 \texttt{vector}의 모든 원소를 반복자로 출력한다. 일반 반복자로 각 원소에 \texttt{3}을 더한 뒤 상수 반복자로 결과를 출력한다.
\end{practicebox}

\begin{lstlisting}
#include <iostream>
#include <vector>

using namespace std;

int main()
{
    vector<int> numbers = {5, 10, 15, 20, 25};

    for (auto it = numbers.begin(); it != numbers.end(); it++)
    {
        cout << *it << " ";
    }

    cout << endl;

    for (auto it = numbers.begin(); it != numbers.end(); it++)
    {
        *it += 3;
    }

    for (auto it = numbers.cbegin(); it != numbers.cend(); it++)
    {
        cout << *it << " ";
    }

    cout << endl;

    return 0;
}
\end{lstlisting}

\subsection{pair}
\begin{practicebox}
이름과 점수를 저장하는 \texttt{pair<string, int>} 세 개를 \texttt{vector}에 넣고 모든 학생의 이름과 점수를 출력한다.
\end{practicebox}

\begin{lstlisting}
#include <iostream>
#include <vector>
#include <string>
#include <utility>

using namespace std;

int main()
{
    vector<pair<string, int>> students;

    students.push_back({"Alice", 95});
    students.push_back({"Bob", 88});
    students.push_back(make_pair(string("Charlie"), 92));

    for (const auto& student : students)
    {
        cout << student.first
             << ": "
             << student.second
             << endl;
    }

    return 0;
}
\end{lstlisting}

\subsection{algorithm}
\begin{practicebox}
정수 \texttt{vector}를 오름차순 정렬한다. 값 \texttt{7}을 찾고, 값 \texttt{3}의 개수를 출력한 뒤 전체 순서를 뒤집는다.
\end{practicebox}

\begin{lstlisting}
#include <iostream>
#include <vector>
#include <algorithm>

using namespace std;

int main()
{
    vector<int> numbers = {7, 3, 9, 3, 1, 7};

    sort(numbers.begin(), numbers.end());

    for (int number : numbers)
    {
        cout << number << " ";
    }

    cout << endl;

    auto it = find(numbers.begin(), numbers.end(), 7);

    if (it != numbers.end())
    {
        cout << "Found: " << *it << endl;
    }
    else
    {
        cout << "Not found" << endl;
    }

    cout << "Count of 3: "
         << count(numbers.begin(), numbers.end(), 3)
         << endl;

    reverse(numbers.begin(), numbers.end());

    for (int number : numbers)
    {
        cout << number << " ";
    }

    cout << endl;

    return 0;
}
\end{lstlisting}

\subsection{lambda}
\begin{practicebox}
학생 이름과 점수를 저장한 \texttt{vector<pair<string, int>>}를 점수 내림차순으로 정렬한다. 점수가 같으면 이름 오름차순으로 정렬한다.
\end{practicebox}

\begin{lstlisting}
#include <iostream>
#include <vector>
#include <string>
#include <utility>
#include <algorithm>

using namespace std;

int main()
{
    vector<pair<string, int>> students = {
        {"Alice", 90},
        {"David", 85},
        {"Bob", 90},
        {"Charlie", 95}
    };

    sort(students.begin(), students.end(),
        [](const pair<string, int>& a,
           const pair<string, int>& b)
        {
            if (a.second != b.second)
            {
                return a.second > b.second;
            }

            return a.first < b.first;
        });

    for (const auto& student : students)
    {
        cout << student.first
             << ": "
             << student.second
             << endl;
    }

    return 0;
}
\end{lstlisting}

\section{실습 파일 구성}
\begin{practicebox}
Day 10 파일은 다음과 같이 관리한다.
\begin{itemize}[leftmargin=1.5em]
    \item \texttt{01\_vector.cpp}
    \item \texttt{02\_iterator.cpp}
    \item \texttt{03\_pair.cpp}
    \item \texttt{04\_algorithm.cpp}
    \item \texttt{05\_lambda.cpp}
    \item \texttt{practice/p01\_vector.cpp}
    \item \texttt{practice/p02\_iterator.cpp}
    \item \texttt{practice/p03\_pair.cpp}
    \item \texttt{practice/p04\_algorithm.cpp}
    \item \texttt{practice/p05\_lambda.cpp}
\end{itemize}
\end{practicebox}

\end{document}
